% A Deterministic Runtime for Agent-Graph Research
%
% 投稿目标:TMLR。正文须 ≤12 页(不含参考文献与附录)——这是把审稿期从 4 周
% 压到 2 周的硬条件,见 jmlr.org/tmlr/reviewer-guide.html。
%
% 本文件刻意不内联 TMLR 样式:投稿前须从 jmlr.org/tmlr/ 下载官方 tmlr.sty
% 并放到本目录。格式违规是 desk reject 的理由,不要用自制样式凑。
%
% 编译:latexmk -pdf main.tex   (本机当前无 LaTeX 工具链,见 README)
%
% ⚠ 所有标为 \TODO 的数字尚未测得。实况实验需要 API 凭据;在跑完之前
%   不得填入任何占位数值——编造的数字即便标注"示意"也会污染后续写作。

\documentclass{article}

% 投稿时替换为:\usepackage{tmlr}
% 匿名评审阶段用 tmlr(默认匿名);camera-ready 用 \usepackage[accepted]{tmlr}
\usepackage[preprint]{neurips_2024}  % 占位,仅供本地排版预览

\usepackage[utf8]{inputenc}
\usepackage[T1]{fontenc}
\usepackage{hyperref}
\usepackage{url}
\usepackage{booktabs}
\usepackage{amsmath}
\usepackage{amsthm}
\usepackage{amssymb}
\usepackage{graphicx}

\newtheorem{theorem}{Theorem}
\newtheorem{definition}{Definition}
\newcommand{\TODO}[1]{\textcolor{red}{[TODO: #1]}}
\usepackage{xcolor}

\title{A Deterministic Runtime for Agent-Graph Research,\\
       and What It Says About Weights and Messages}

\author{Anonymous}

\begin{document}
\maketitle

\begin{abstract}
Comparisons between multi-agent LLM topologies are difficult to reproduce.
Wall-clock concurrency, per-arm differences in budget accounting, and divergent
partial-failure handling all perturb the outcome, so a reported topology effect
is confounded with runtime nondeterminism and with unequal effective budgets.
We present an agent-graph runtime built so that \emph{the runtime contributes
zero entropy}: given a recorded sequence of model responses, the final run
snapshot is reproducible bit-exactly regardless of concurrency, scheduling, or
completion order. Three properties make this possible --- order-invariant
super-step merging, fail-closed all-or-nothing ingestion of agent proposals, and
a budget ledger that enforces parity by construction rather than by post-hoc
filtering. Treating the runtime as a measurement instrument, we then ask what
agent-graph mechanisms actually buy at matched token budget, ablating weighted
scheduling, budget allocation, cross-generation skip context, and message
passing as independent toggles with everything else held bit-identical.
\TODO{one-sentence summary of the empirical findings}
All experiments ship with replay transcripts, so every number can be re-derived
offline at zero API cost.
\end{abstract}

\section{Introduction}
\label{sec:intro}

\TODO{1.25pp. The reproducibility gap in multi-agent LLM evaluation. Three
claims: (i) conditional determinism --- the runtime contributes zero entropy;
(ii) a fail-closed trust boundary in which agents propose and only the
controller mutates the graph; (iii) budget parity by construction. Then preview
the mechanism findings.}

The claim in (i) needs stating precisely, because the obvious stronger version
is false. At temperature $>0$ two live runs of the same configuration do not
agree, and we do not claim they do. What we claim and verify is that the
\emph{runtime} is not a source of variation: fixing the model's responses fixes
the final snapshot digest exactly, under any concurrency limit and any
completion order. All remaining entropy in an experiment is therefore
attributable to the model, which is what makes a topology comparison meaningful.

\section{Related work}
\label{sec:related}

\TODO{1.0pp. Three paragraphs plus one compressed capability table.}

Recursive agent runtimes --- LangGraph~\citep{langgraph}, AutoGen's
GraphFlow~\citep{wu2023autogen}, MetaGPT~\citep{hong2024metagpt},
CrewAI~\citep{crewai}, CAMEL~\citep{li2023camel}, and
Swarm~\citep{openaiswarm} --- organise agents as trees or predefined workflows
and offer no learnable topology. Graph-optimisation work ---
GPTSwarm~\citep{zhuge2024gptswarm}, DyLAN~\citep{liu2024dylan},
AgentPrune~\citep{zhang2024agentprune}, G-Designer~\citep{zhang2024gdesigner},
AutoAgents~\citep{chen2023autoagents}, and MacNet~\citep{qian2024macnet} ---
treats the graph as optimisable but fixes it before execution. Neither camp
offers determinism guarantees or budget parity across compared configurations,
which is the gap this work addresses.

\section{Runtime model and determinism}
\label{sec:runtime}

\TODO{2.5pp.}

\subsection{Attributed graph over time}

\TODO{Formalise $G_t = (V_t, E_t, X_t, W_t)$; four typed edge kinds and why a
single edge type is insufficient; cycle unrolling to
$(\mathrm{agent\_id}, \mathrm{round}, \mathrm{invocation\_id})$.}

\subsection{Proposals and the trust boundary}

Agents cannot mutate the graph. A model's \texttt{graph\_\_*} tool calls produce
only \emph{proposals} in the returned result; the controller re-validates every
proposal at ingestion and materialises nodes and edges itself. Ingestion is
fail-closed and all-or-nothing: an invalid result adopts \emph{none} of its
proposals and commits its entire token reservation.

\subsection{Super-step execution}

Execution proceeds in BSP super-steps~\citep{valiant1990bsp} with three phases:
freeze the ready frontier and claim reservations under a lock; execute claimed
runners concurrently with \emph{no} lock held; re-acquire the lock and merge
outcomes \emph{ordered by claim sequence rather than completion order}.

\begin{theorem}[Order-invariant merge]
\label{thm:order}
Fix a root specification, a controller policy, and the multiset of agent results
produced in a run. The final snapshot is independent of the order in which
runners complete and of the concurrency limit.
\end{theorem}

\TODO{Proof sketch: merge iterates outcomes sorted by \texttt{execution.sequence},
a value assigned inside the critical section; all records are immutable and
replaced rather than mutated in place; all mutators are deterministic functions
of (state, result) with identifiers derived by hashing. Hence the merge is a
fixed fold over a fixed sequence.}

\begin{theorem}[Budget parity]
\label{thm:budget}
On every termination path --- success, structured failure, invalid result,
runner exception, cancellation, and dependency block --- committed tokens never
exceed the configured limit, and reserved tokens return to zero at convergence.
\end{theorem}

\TODO{Proof sketch: enumerate the six paths against the ledger transitions
(reserve / claim / commit-success / commit-full / release).}

\subsection{Audit chain and resumability}

\TODO{SHA-256 chains store digests only, never plaintext. Checkpoint v1 is pure
snapshot replay; v2 adds a resumable slice so completed nodes replay from an
archive without re-running or re-billing.}

\section{Mechanisms under study}
\label{sec:mechanisms}

\TODO{2.0pp.}

\subsection{Four disjoint scalars}

\begin{description}
\item[\texttt{priority}] orders the ready frontier and nothing else.
\item[\texttt{budget\_share}] sizes a child's token allocation.
\item[\texttt{context\_gate}] bounds how much of an edge's payload may enter the
      target's context.
\item[\texttt{trust}] gates and ranks delivery.
\end{description}

Keeping these disjoint is deliberate: a single fused ``weight'' conflates
scheduling, economics, information flow, and credibility, after which no
ablation can attribute an effect to any one of them.

\paragraph{Why trust learning must be cross-run and role-keyed.}
This is a structural consequence worth stating as a result rather than a
limitation. In a runtime where a node executes exactly once and never re-enters,
within-run node-level trust learning is impossible by construction: for node
$n$'s trust to influence any decision it would have to be updated after $n$
produces its result and then read again, but $n$ never runs again, and its
children froze their identities at spawn time. The only feasible learning is
therefore across runs and aggregated by role. We use an interpretable
exponential moving average, $\mathrm{trust}^{(t+1)} = (1-\alpha)\,
\mathrm{trust}^{(t)} + \alpha\, r^{(t)}$, and explicitly decline end-to-end
reinforcement learning at this stage; credit assignment is by leave-one-out
role ablation.

\subsection{Typed context delivery}

\TODO{Message injection, \texttt{depends\_on} result propagation (fan-in), skip
gating, and the token knapsack. Explain the mandatory/discretionary split:
\texttt{depends\_on} is declarative data flow that a lexical relevance score has
no authority to veto, whereas messages and skip context are discretionary
evidence subject to routing.}

Routing scores a candidate delivery as
$s_{u \rightarrow v} = \mathrm{priority}_u \cdot \mathrm{rel}(u,v) \cdot
\mathrm{trust}_{\mathrm{role}(u)}$, selects Top-$K$, then fills a token budget
greedily. Relevance is a deterministic lexical IDF-weighted cosine. Embeddings
are excluded on purpose: they would add a runtime dependency and, if served
remotely, a source of nondeterminism at the centre of the zero-entropy claim.
We therefore report lexical relevance as a \emph{lower bound} on what routing
can achieve, which makes a positive result conservative and obliges us to hedge
a null one.

\subsection{Why black-box skip connections are semantic, not residual}

\TODO{With black-box APIs there is no access to hidden states, so
$h_{l+1} = F(h_l) + h_l$ is unavailable. What is achievable is a semantic skip:
relevance filter, instruction isolation, structured summary, gating, and
injection as evidence rather than as a system instruction.}

\section{Experimental design}
\label{sec:design}

\TODO{1.5pp.}

\subsection{Arms and toggles}

\begin{table}[t]
\centering
\caption{Experimental arms. All arms share one token ceiling, enforced by the
controller's ledger rather than by discarding overspent runs.}
\label{tab:arms}
\begin{tabular}{llp{7.2cm}}
\toprule
Arm & Label & Increment \\
\midrule
A1 & single agent & spawn tool withheld; \texttt{maxNodes}~$=1$ \\
A2 & equal-weight tree & $+$ spawn, uniform \texttt{budget\_share}, fixed \texttt{priority} \\
A3 & weighted tree & $+$ model-chosen \texttt{priority} and \texttt{budget\_share} \\
A4 & fixed skip graph & $+$ cross-generation skip context \\
A5 & dynamic graph & $+$ messages and role-keyed trust \\
\bottomrule
\end{tabular}
\end{table}

Two orthogonal toggles, \textsc{fan-in} and \textsc{delivery}, ablate the
mechanisms behind the communication claim. The delivery toggle is implemented by
closing the context-injection token budget rather than by branching, so control
and treatment traverse identical code paths and differ only in available budget.

\subsection{Tasks, scoring, and budget parity}

Scorers are deterministic pure functions. This is load-bearing rather than
fastidious: variance is a reported metric, so a stochastic scorer would mix
scorer variance into the quantity we attribute to the model. That requirement is
also the honest reason only two task families survive --- unit-testable code
generation~\citep{chen2021humaneval,liu2023evalplus} and mathematics with
deterministic answers~\citep{cobbe2021gsm8k}.

\subsection{Objective and the weight sweep}

We report $R = \mathrm{Quality} - \lambda\,\mathrm{Cost} - \mu\,\mathrm{Latency}
- \nu\,\mathrm{Failure}$, but treat the Pareto front over quality and cost as
the primary figure. The reason is empirical, not stylistic: in pilot runs the
ranking of arms under $R$ flips as $\lambda$ crosses a threshold, so any single
choice of $\lambda$ would be an unjustifiable degree of freedom. We therefore
sweep $\lambda/\mu/\nu$ and report whether conclusions survive the sweep.

\subsection{Tiered compute and what is not covered}

\begin{table}[t]
\centering
\caption{Tiered experiment matrix. The frontier tier is a deliberate
confirmatory subset, reported here rather than left for a reader to infer from
row counts.}
\label{tab:matrix}
\begin{tabular}{llll}
\toprule
Tier & Conditions & Seeds & Cells \\
\midrule
small    & 7 (5 arms $+$ 2 ablations) & 5 & \TODO{} \\
mid      & 7                          & 5 & \TODO{} \\
frontier & 3 (A1, A3, A5)             & 3 & \TODO{} \\
\bottomrule
\end{tabular}
\end{table}

\section{Results}
\label{sec:results}

\TODO{2.75pp. No numbers may be written here until the live experiments have
run. Each subsection below states what will be reported.}

\subsection{Determinism validation}

\TODO{Replay digest match rate --- must be 100\%; order-shuffle invariance
across concurrency limits and completion orders. A rate below 100\% would mean
an undiscovered source of nondeterminism and would invalidate
Theorem~\ref{thm:order}; we report it as a stop-the-line check rather than as a
summary statistic.}

\subsection{Weights}

\TODO{priority vs budget\_share vs trust, ablated. Paired per-instance
comparisons with confidence intervals, not bare mean differences.}

\subsection{Messages and fan-in}

\TODO{delivery on/off, fan-in on/off, skip on/off; skip-utilisation rate.}

\subsection{Robustness to an injected faulty agent}

\TODO{Contamination of the final answer by a confidently-wrong role, with and
without the minimum-trust gate.}

\subsection{Cross-model stability}

\TODO{Whether effects replicate across model scales. An effect that appears on
one model only is reported as model-specific --- a finding, not a failure.}

\section{Limitations and conclusion}
\label{sec:limits}

Determinism is conditional on the model: we claim zero runtime entropy, not
bit-identical live runs. Relevance is lexical rather than semantic, so routing
results are a lower bound. Trust learning is cross-run and role-level, which
follows from the single-shot node model and cannot be loosened without changing
that model. Two task families are covered; the four families without
deterministic judges are discussed in the appendix but not measured. No learned
topology optimiser is included --- rule-based and statistical mechanisms only.
Finally, the code-generation scorer executes model-generated code under
process-level isolation (separate subprocess, blocked network and
subprocess modules, wall-clock and memory limits, whitelisted environment,
disposable working directory). This is process isolation, not a sandbox: the
child runs as the invoking user. We state this plainly in the artifact rather
than describing the setup as sandboxed.

\TODO{Conclusion paragraph once results exist.}

\bibliographystyle{plainnat}
\bibliography{refs}

\appendix

\section{Framework survey}
\TODO{Full 14-framework comparison across seven capability axes.}

\section{Rejection codes}
\TODO{The validation codes enforced at the ingestion boundary.}

\section{Hash chain and checkpoint schemas}
\TODO{Event chain schema; checkpoint v1 and v2 formats.}

\section{Per-instance results and cost tables}
\TODO{}

\section{Mechanisms considered and declined}
\TODO{Learned topology optimiser; hidden-state residual connections;
end-to-end RL credit assignment. Include the feasibility scores and the
reasoning for declining each.}

\end{document}
